\documentclass[11pt]{article}
\usepackage{fontspec}
\setmainfont{Times New Roman}
\setsansfont{Arial}
\setmonofont{Consolas}
\usepackage{amsmath,amssymb,mathtools}
\usepackage{geometry}
\usepackage{microtype}
\geometry{a4paper,margin=2.35cm}
\setlength{\parindent}{1.25em}
\setlength{\parskip}{0pt}
\makeatletter
\renewcommand\footnotesize{\@setfontsize\footnotesize{7.7}{8.7}\selectfont}
\makeatother
\emergencystretch=95em
\allowdisplaybreaks
\sloppy
\hbadness=10000
\vbadness=10000
\hfuzz=2pt
\vfuzz=10pt
\raggedbottom

\begin{document}

% Unit: Paper31 Section03 entry 3 v001. Direct Interslavic Latin translation from the German final-audited source slice, source lines 198--200 / segments P31-S0044--P31-S0045. Footnote counter continues after Paper31 Section03 entry 2.
\setcounter{footnote}{22}

2b. Iz polnoj razložimosti prvogo roda kolca $\mathfrak R$ slěduje polna razložimost prvogo roda kolca $\overline{\mathfrak R}$; občneje, polna razložimost prvogo roda vsakogo medžukolca medžu $\mathfrak R$ i $\overline{\mathfrak R}$.

Poněže po \S 2, 3 predstavjenju $\mathfrak R$ kako direktnoj sumy odgovarja tako predstavjenje $\overline{\mathfrak R}$, že vsaka komponenta $\overline{\mathfrak R}_i$ stava razširjeno kolco $\mathfrak R_i[\overline P_i]$ --- gde $\overline P_i$ jest razširjeno polje izomorfno k $\overline P$ od podpolja $P_i$ komponenty $\mathfrak R_i$, izomorfnogo k $P$, --- dostatočno jest razsmotriti pojedinu komponentu $\mathfrak R_i$. Zato bez ograničenja obćnosti možno predpolagati, že $\mathfrak R$ jest polje, imeno konečno razširjeno polje prvogo roda svojego podpolja $P$. Tedy $\mathfrak R$ nastavaje prisojedinjenjem primitivnogo elementa prvogo roda k $P$; drugymi slovami, $\mathfrak R$ jest izomorfno kolcu klasov ostatkov $P[x]/f(x)$, gde $f(x)$ označaje nerazložimy polinom prvogo roda v polinomnoj oblasti $P[x]$ jednej neopreděljenej $x$. Ako $Q$ jest razširjeno polje od $P$, togda $Q[x]/f(x)$ jest izomorfno k $\mathfrak R[Q]$; bo pri prěhodu od $P[x]$ k $Q[x]$ rang kolca klasov ostatkov sohranja se i jest ravny stepeni $f(x)$, takže ide točno o konstrukciju razširjenogo kolca danu v \S 1, 2. V $Q[x]$ polinom $f(x)$ razpada se na po parah vzajemno proste nerazložime polinomy prvogo roda, od kojih vsaky porađaje prosty ideal v $Q[x]$; to znači, že $Q[x]/f(x)$, a zaradi izomorfizma takože $\mathfrak R[Q]$, stava polno razložimo prvogo roda. Poněže to važi za vsako razširjeno polje $Q$, osoblivo takože za algebraično zatvorjeno polje $\overline P$, s tym jest 2b dokazano.

\end{document}
